\chapter{Implementation and Methodology}

\section{Implementation Goals}

The implementation part of this thesis was designed as an experimental and
educational tool for studying the Merkle--Hellman knapsack cryptosystem. Its
main goal is to show the full workflow of the scheme. This includes key
generation, encryption, decryption, and checking that the recovered plaintext
is correct. Besides the classic variant it also includes a permuted variant and
a simple iterated variant.

The program is not intended to be a secure cryptographic library. The schemes
studied in this thesis are historically important, but they are known to be
insecure against cryptanalytic attacks. The implementation is used mainly to
illustrate how the trapdoor construction works in practice and to support the
theoretical discussion from the previous chapters.

The implementation also includes simple solvers for small subset-sum instances.
They are used to illustrate how direct recovery of the message becomes harder
as the message length increases. Our intention was not to implement Shamir's
structural attack or lattice-based attacks.

Another goal was to prepare a simple environment for experiments. The program
supports deterministic runs through explicit seeds, so the same keys and test
cases can be reproduced across multiple runs. This is useful for debugging,
correctness testing, and benchmarks of running time, density, and size-related
key parameters such as the bit length of the private-weight sum.

\section{Choice of Programming Language and Libraries}

The implementation is written in C. This choice was motivated by both practical
and educational reasons. C is commonly used in systems programming and in many
cryptographic libraries, where performance, memory layout, and low-level
arithmetic are important. For this thesis, using C also made it possible to work
closer to the representation of keys, bit vectors, and large integer values
instead of relying on a high-level prototype. It was also useful as practice
with manual memory management, explicit error handling, and build tooling.

A major requirement of the implementation is support for integers larger than
those provided by the standard C integer types. In knapsack cryptosystems, the
weights, modulus, multiplier, and ciphertexts can grow quickly as the block
length increases. For this reason, the implementation uses the GNU Multiple
Precision Arithmetic Library (GMP). GMP is used for large integer arithmetic,
such as additions, comparisons, modular inversion, and ciphertext computation.
This made it possible to keep the implementation in C while supporting
arithmetic beyond standard 64-bit limits.

The random generator used by the program is based on PCG~\cite{ONeill2014PCG}.
It is used for the randomized parts of the implementation, such as generating
private weights, permutations, and test messages. The generator is suitable for
experiments and demonstrations, but the program does not claim to provide
production cryptographic randomness.

The build process is controlled by a simple Makefile. It supports separate
debug and release builds, as well as an address-sanitized configuration for
runtime checking. This was useful during development, because the code uses
dynamic allocation for keys, messages, ciphertexts, and scheme-specific
internal structures. The toolchain was kept simple so that the program could be
built and tested in both macOS and Linux environments with only a C compiler,
Make, and GMP installed.

\section{Program Architecture}

The implementation was organized as a small modular C program. The main goal
was to keep the cryptographic code separate from user interaction, input
parsing, benchmarking, and subset-sum solvers. This structure makes the
implementation easier to follow and makes later extensions, such as new modes or
scheme variants, less dependent on changes in unrelated code.

The executable entry point in \texttt{main.c} is intentionally small. It
hands control to the application layer in \texttt{app.c}, which selects the
requested mode and coordinates the rest of the program. Command-line parsing
and validation are handled in \texttt{cli.c}. This keeps input handling
separate from the implementation of the schemes. Common success and error
states are represented by the \texttt{KnapStatus} type, which is used across
the program for explicit error handling.

The program supports demo mode, benchmark mode, and a solver mode.
Demo mode can encrypt either a raw bit string or a plaintext message. Plaintext
input is converted to bits, split into fixed-size blocks, encrypted block by
block, and then reconstructed after decryption. Benchmark mode repeats key
generation, encryption, and decryption and records timing data.
In the command-line interface, the solver mode is named \texttt{attack}, but it
only runs brute-force search and a meet-in-the-middle solver on small generated
subset-sum instances. In the current program, this mode is intended for the
classical Merkle--Hellman scheme.

\subsection{Scheme Interface and Variants}

The common scheme interface is a small abstraction defined in
\texttt{scheme.h}. Each variant is then represented by a \texttt{SchemeOps}
structure containing function pointers for key generation, encryption,
decryption, and cleanup, together with a short identifier and a display name.
This gives all variants the same external interface, while still allowing each
one to keep its own internal logic and data structures. 

\begin{lstlisting}[language=C,caption={Common interface.},label={lst:scheme-ops}]
    typedef struct {
        SchemeInfo info;
        KnapStatus (*keygen)(const SchemeKeygenParams *params,
                            SchemeKey *out_keypair);
        KnapStatus (*encrypt)(const SchemeKey *keypair, 
                            BitView message,
                            mpz_t out_ciphertext);
        KnapStatus (*decrypt)(const SchemeKey *keypair,
                            const mpz_t ciphertext,
                            BitBuf *out_message);
        void (*scheme_key_clear)(SchemeKey *keypair);
    } SchemeOps;
\end{lstlisting}


The implemented variants are the classical Merkle--Hellman scheme, a permuted
variant, and a simple iterated variant. Their code is placed in separate source
files. Scheme selection is handled in \texttt{scheme.c}, where the identifier
passed on the command line is mapped to the corresponding implementation. The
identifier \texttt{mh} is also accepted as an alias for \texttt{mh-classic}.

The classical variant implements the standard Merkle--Hellman scheme with a
superincreasing private sequence, a modulus, and a multiplier. During key
generation, the private weights are created by maintaining the current sum
\(S\) of previous private weights. For each new weight, a random increment
\[
    \delta_i \in \{1,\dots,\texttt{MH\_DEFAULT\_DELTA\_MAX}\}
\]
is chosen and the next private weight is set to
\[
    w_i = S + \delta_i.
\]
After all private weights are generated, the modulus is chosen as the total
private sum plus a small random margin.
\[
    m = S + \mu,
\]
where
\[
    \mu \in \{1,\dots,\texttt{MH\_DEFAULT\_MARGIN\_FACTOR}\cdot n\}.
\]
The multiplier is then sampled from the range \(2 \le r < m\) until it has
an inverse modulo \(m\). These constants are later varied in the
parameter-sweep experiments.

The permuted variant adds a secret permutation of the private weights before
the public weights are created. The iterated variant applies more than one
modular transformation layer and removes these layers in reverse order during
decryption. In the command-line interface, the default iterated variant uses two
modular transformation layers. The option \texttt{--layers} can be used with
\texttt{mh-iterated} to choose a value from 2 to 20 layers.

\subsection{Data Representation and Support Modules}

The schemes work with fixed-length bit blocks. Raw bit-string input can be used
directly, while plaintext input is first converted from bytes to bits. The bit
sequence is then split into blocks of the selected size, and each block is
encrypted as one knapsack ciphertext. For plaintext demo input, the default
block size is 128 bits unless another value is given with the \texttt{--n}
command-line option.

The implementation uses \texttt{BitBuf} for owned bit arrays and
\texttt{BitView} for non-owning views into an existing bit array. This is used
when passing individual blocks to encryption. Plaintext strings are stored in
\texttt{TextBuf}, but the schemes themselves only work with bits.

Block splitting is handled by \texttt{BitBlocks}. It stores the padded bit
buffer, the block size, the number of blocks, and the original bit length. The
buffer is initialized with zeros before the original bits are copied into it, so
unused positions in the last block become zero padding. After decryption, only
the original number of bits is used when printing the recovered bits or
converting them back to text.

Several support modules are separated from the scheme implementations. Seed
handling is implemented in \texttt{seed.c} and supports runs with an explicit
seed or with entropy from the operating system. This seed handling is used by demo, benchmark, and solver
runs, so generated keys and messages can be repeated when needed.

Benchmarking is implemented in \texttt{bench.c}. It measures key generation,
encryption, decryption, and total running time, and records values such as
density, margin, and the bit length of the private-weight sum.

The subset-sum solvers are implemented in \texttt{attack\_mh.c}. Output
formatting is separated into \texttt{output.c}, and plaintext conversion is
handled in \texttt{plaintext.c}.

\section{Testing and Experiment Support}

Correctness was checked with the script \texttt{tests/test.sh}. The script
builds the address-sanitized version of the program and runs selected demo,
benchmark, and solver cases. It checks that the implemented variants can
encrypt and decrypt raw bit strings and plaintext messages. It also checks that
the solver mode can recover generated messages for small parameters.

The tests are not a formal verification of the schemes. They are used to check
that the implemented program behaves as expected for the tested inputs and that
basic error cases are handled. This is sufficient for the purpose of this
thesis, because the implementation is used as an educational and experimental
tool.

We used Python outside the C program for experiment automation and plotting.
The scripts in the \texttt{experiments} directory do not implement the
cryptographic schemes themselves. Instead, they run the compiled C executable
with selected parameters, read the CSV output produced by the program, store the
results, and generate the plots used in the results chapter.

The automation scripts use only simple Python tooling. The standard library
modules \texttt{pathlib}, \texttt{subprocess}, \texttt{csv}, and \texttt{re}
are used for locating files, executing the compiled program, parsing CSV output,
and modifying selected compile-time constants during parameter sweeps. The
plotting script uses \texttt{pandas} for grouping and averaging the measured CSV
data and \texttt{matplotlib} for creating the figures.

This separation is important for interpreting the measurements. Key generation,
encryption, decryption, and the subset-sum solvers are executed by the C
implementation. Python is used only as an experiment and visualization
tool. 
